\section{L'accessibilité en ligne et les accommodations éducatifs}\label{laccessibilite-en-ligne-et-les-accommodations-educatifs}

\stanceinfo{\textbf{FCEG CCLI 2024} [2024-01-05]}{2024-01-05}

\officialstance{La Fédération canadienne étudiante de génie estime que les ressources pédagogiques [enregistrements, notes de cours, etc.] devraient rester accessibles en ligne au même niveau de qualité que les sessions en personne. La FCEG estime en outre que des aménagements doivent être mis en place pour les évaluations tout au long du semestre universitaire en cas de circonstances atténuantes, de maladie ou d'autres problèmes personnels, afin d'offrir un soutien adapté à tous les étudiants en génie.}

\stanceheading{La position des étudiants}
\begin{itemize}
 \item Tous les étudiants en ingénierie devraient avoir les mêmes chances de poursuivre et d'exceller dans leur formation, quelles que soient leurs capacités et circonstances individuelles.
 \item La COVID-19 a offert la possibilité de disposer de ressources en ligne permanentes, apprécié par les étudiants en ingénierie pour améliorer leurs performances académiques.
 \item Une approche de conception universelle des ressources éducatives et des examens ne tient pas compte de la diversité des besoins et des difficultés d'apprentissage.
 \item Les évaluations traditionnelles ne reflètent pas la compréhension du contenu dans toutes les situations en raison de l'environnement très stressant et de la nécessité d'être performant sur le champ.
 \item Des aménagements devraient être prévus en fonction des besoins afin de répondre aux besoins des étudiants en ingénierie, de les soutenir dans leur formation et d'encourager leur réussite.
\end{itemize}

\stanceheading{Les recommandations aux partenaires, aux parties intéressées et aux autres entités}
\begin{itemize}
 \item La FCEG demande au BCAPG d'encourager les établissements dans le processus de mise à jour de leurs politiques institutionnelles et d'être plus inclusifs à l'égard de l'enseignement hybride et du matériel permanent en ligne.
 \item La FCEG demande au BCAPG d'évaluer les aménagements et l'accessibilité disponibles dans un programme d'ingénierie lors d'une visite d'accréditation.
 \item La FCEG demande aux DDIC d'encourager les établissements à créer un modèle hybride standard pour la prestation des cours en fonction des commentaires des étudiants sur le programme.
 \item La FCEG demande aux DDIC de veiller à ce que des accommodements soient accordés aux personnes ayant des circonstances ou des désavantages divers afin de promouvoir l'équité dans le domaine de l'ingénierie.
\end{itemize}

\stanceheading{Ce que fait la FCEG}
\begin{itemize}
 \item L'enquête nationale de 2023 a recueilli des informations sur les ressources académiques et leur disponibilité.
 \item Des analyses documentaires sur l'éducation dans un monde post-pandémique qui ont trouvé des avantages à l'apprentissage en ligne.
 \item Le FCEG continue de rencontrer les parties prenantes au sujet de la flexibilité de l'éducation et de la qualité de l'enseignement.
\end{itemize}

\stanceheading{Ce que la FCEG envisage de faire}
\begin{itemize}
 \item La FCEG travaillera avec les parties prenantes et les partenaires pour trouver un terrain d'entente afin d'entamer un travail politique sur les aménagements en ligne et l'accessibilité académique dans les formations d'ingénieurs.
 \item La FCEG créera des ressources pour aider ses membres à défendre l'accessibilité en ligne et les aménagements académiques par le biais de lettres de soutien.
 \item La FCEG recueillera davantage d'informations et de témoignages sur les politiques et les expériences actuelles des étudiants en matière d'accessibilité en ligne et d'aménagements académiques.
\end{itemize}

\newpage
